\documentclass[a4paper]{article}

\usepackage{graphicx}
\usepackage{amsmath,amssymb}
\usepackage{minted}
\usepackage{multirow}
\usepackage{float}
\usepackage{todonotes}
\usepackage{forest}
\usepackage{xstring}
\usepackage{xcolor}
\usepackage[most]{tcolorbox}
\usepackage{xparse}
\usepackage{natbib}

\usetikzlibrary{graphs,graphdrawing}
\usegdlibrary{layered}

% for external graphviz
\input{/home/donkarlo/Dropbox/repo/nd_language_ai_project/src/nd_language_ai/formal/computer/programming/tex/latex/code_clips/external_graphviz.tex}

% for tags
\input{/home/donkarlo/Dropbox/repo/nd_language_ai_project/src/nd_language_ai/formal/computer/programming/tex/latex/parts/control_sequence/kind/semtag/semtag.tex}

% for links
\input{/home/donkarlo/Dropbox/repo/nd_language_ai_project/src/nd_language_ai/formal/computer/programming/tex/latex/code_clips/seehere.tex}

\title{Seit, ab und von an}
\date{}

\begin{document}
    \maketitle
    
    
    \section{seit}
        \textit{seit} verwendet man, wenn etwas in der Vergangenheit begonnen hat und bis jetzt andauert.
        
        \begin{itemize}
            \item Ich wohne seit einem Jahr in Graz.
            \item Er lernt seit zwei Stunden Deutsch.
        \end{itemize}
    
    
    \section{ab}
        \textit{ab} verwendet man, wenn etwas von einem bestimmten Zeitpunkt an beginnt. Dieser Zeitpunkt liegt meistens in der Gegenwart oder in der Zukunft.
        
        \begin{itemize}
            \item Ab morgen lerne ich jeden Tag Deutsch.
            \item Ab nächster Woche arbeite ich mehr.
        \end{itemize}
    
    
    \section{von ... an}
        \textit{von ... an} hat eine ähnliche Bedeutung wie \textit{ab}, betont aber den Anfangspunkt stärker. Es bedeutet: von diesem Zeitpunkt an und danach.
        
        \begin{itemize}
            \item Von morgen an lerne ich jeden Tag Deutsch.
            \item Von diesem Tag an war alles anders.
        \end{itemize}
    
    
    \section{Unterschied zwischen seit und ab}
        \textit{seit} beschreibt einen Beginn in der Vergangenheit mit Wirkung bis in die Gegenwart.
        \textit{ab} beschreibt einen Beginn von einem bestimmten Zeitpunkt an.
        
        \begin{itemize}
            \item Ich wohne seit einem Jahr in Graz.
            \item Ab nächstem Monat wohne ich in Graz.
        \end{itemize}
    
    
    \section{Unterschied zwischen ab und von ... an}
        Beide Ausdrücke bezeichnen einen Beginn ab einem bestimmten Zeitpunkt.
        \textit{von ... an} hebt diesen Zeitpunkt meistens stärker hervor.
        
        \begin{itemize}
            \item Ab morgen mache ich mehr Sport.
            \item Von morgen an mache ich mehr Sport.
        \end{itemize}
    
    
    \section{Kurze Zusammenfassung}
        \begin{itemize}
            \item \textit{seit}: Beginn in der Vergangenheit, andauernd bis jetzt
            \item \textit{ab}: Beginn von einem bestimmten Zeitpunkt an
            \item \textit{von ... an}: Beginn von einem bestimmten Zeitpunkt an, mit stärkerer Betonung des Anfangspunkts
        \end{itemize}
    
    %\bibliographystyle{plainnat}
    %\bibliography{/home/donkarlo/Dropbox/repo/nd_shared_research/bibliography/bibliography.bib}
\end{document}